\chapter{MACSLAM: Unified Cross-Modal Uncertainty-Aware Visual-Inertial SLAM} \label{chapter:macslam}
\label{chapter:macvio}
\fancyhead[LO]{\makebox[\textwidth][r]{Chapter 8. MACSLAM}}

\section{Introduction}

This chapter presents MACSLAM, a unified SLAM framework that merges MACVIO-style cross-modal visual-inertial fusion with multi-frame optimization, loop closure, and semantic mapping under one uncertainty-aware formulation.
It extends modules from earlier chapters: metrics-aware visual covariance from MAC-VO, inertial uncertainty modeling from AirIMU and AirIO, and robust initialization and online calibration from MAC-I$^2$.
The local odometry layer learns cross-modal uncertainty correlations and converts them into adaptive fusion weights, while the global SLAM layer propagates those learned covariances into factor-graph optimization.

The key design choice is a unified visual foundation front-end based on DINO features.
The same visual representation is used for local geometric matching, global visual place recognition, and semantic-language alignment through lightweight adapters.
This design aims to reduce front-end fragmentation and make local odometry, loop closure, and semantic mapping mutually consistent.

The motivation for MACSLAM is that robust autonomy requires local fusion and global consistency to be solved together.
A deployable system must adapt visual-inertial trust under short-term degradation, maintain global consistency over long trajectories, recover from drift through loop closure, and support higher-level map understanding for planning and interaction.
Existing pipelines often treat these capabilities as loosely connected modules with separate feature spaces, which introduces inconsistency and brittle behavior under domain shift.

MACSLAM addresses this fragmentation through three unification principles.
First, visual and inertial information are fused through a learned reliability state that determines how much each modality should influence estimation at each time step.
Second, the same uncertainty-aware fusion signal is used inside a multi-frame factor graph rather than being confined to local odometry.
Third, one foundation visual representation supports local matching, global place recognition, and semantic-language mapping.
This shared representation is intended to improve both robustness and scalability while opening a path toward language-aware semantic SLAM in real deployments.

\section{Method}
\label{sec:macslam:method}

\subsection{Unified Optimization Objective}

Let $\mathcal{X}_{1:T}$ be the trajectory states, $\mathcal{M}$ the map variables, and $\mathcal{L}$ loop-closure constraints.
MACSLAM solves a single factor-graph objective:
\begin{equation}
\label{eq:macslam:objective}
(\mathcal{X}_{1:T}^{*}, \mathcal{M}^{*})
=
\argmin_{\mathcal{X}_{1:T},\mathcal{M}}
\sum_{k\in\mathcal{V}} w_k^{v}\, \left\| r_k^{v} \right\|_{{\Sigma}_k^{v}}^{2}
+
\sum_{k\in\mathcal{I}} w_k^{i}\, \left\| r_k^{i} \right\|_{{\Sigma}_k^{i}}^{2}
+
\sum_{k\in\mathcal{L}} \left\| r_k^{lc} \right\|_{{\Sigma}_k^{lc}}^{2},
\end{equation}
where visual, inertial, and loop-closure residuals are jointly optimized with learned, metrics-aware covariance.
The adaptive weights $(w_k^{v}, w_k^{i})$ transfer the MACVIO local fusion mechanism into the global SLAM graph.
Unlike conventional pipelines with fixed noise assumptions, MACSLAM propagates data-conditioned uncertainty through all factor types.

\subsection{Cross-Modal Reliability Learning}

MACSLAM uses the MACVIO fusion layer to treat visual and inertial uncertainty as complementary signals.
Given visual measurements $\mathcal{Z}_{1:T}^{v}$ and inertial measurements $\mathcal{Z}_{1:T}^{i}$, the system learns a cross-modal reliability state from uncertainty and context:
\begin{equation}
\label{eq:macslam:reliability_state}
\mathbf{s}_k = \phi\!\left({\Sigma}_k^{v}, {\Sigma}_k^{i}, \mathbf{c}_k\right),
\end{equation}
where $\mathbf{c}_k$ includes motion and feature-quality context such as image degradation, feature density, inertial excitation level, vibration, and bias drift.
The reliability state is mapped to adaptive fusion weights by:
\begin{equation}
\label{eq:macslam:fusion_weights}
\left(w_k^{v}, w_k^{i}\right) = g(\mathbf{s}_k), \quad w_k^{v}, w_k^{i}\in[0,1].
\end{equation}
This mechanism increases trust in the currently reliable modality and suppresses degraded measurements before they can destabilize the local trajectory or the global map.

\subsection{From Local MACVIO to Global MACSLAM}

The merged system keeps the MACVIO role as the local visual-inertial estimator and embeds it inside the MACSLAM graph.
MACVIO provides adaptive visual-inertial residual weights for short-horizon odometry, while MACSLAM reuses the same learned reliability signal for keyframe selection, loop-factor weighting, and global graph consistency.
By combining these layers, the system avoids a common failure mode in which local VIO rejects or over-trusts a modality while the global SLAM back-end assumes fixed measurement quality.

\subsection{DINO-Based Unified Visual Front-End}

MACSLAM uses DINO features as a shared visual backbone with three heads:
\begin{itemize}
    \item \textbf{Geometric head:} dense feature matching and correspondence confidence for local visual odometry factors.
    \item \textbf{Global retrieval head:} place descriptors for loop-closure candidate retrieval and verification.
    \item \textbf{Semantic-language head:} adapter from DINO tokens to semantic and language-aligned embeddings for semantic mapping.
\end{itemize}

The core hypothesis is that one foundation representation can support both local geometry and global semantics more consistently than separate task-specific front-ends.
This allows loop closure and semantic mapping to directly reuse features already optimized for odometry.

\subsection{Visual-Inertial Fusion in Multi-Frame SLAM}

Visual and inertial factors are fused under uncertainty-aware weighting.
Visual covariance priors come from the metrics-aware modeling of MAC-VO, while inertial covariance priors come from AirIMU/AirIO and are adapted online through cross-modal reliability control.
Fusion is not implemented as static modality weighting; instead, confidence is adjusted online according to degradation context (for example, motion blur, low texture, aggressive maneuvers, and IMU disturbance).

\subsection{Loop Closure and Semantic Mapping}

Loop closure is modeled as a two-stage process: DINO-based global retrieval proposes candidates, and geometric verification adds loop factors with uncertainty-aware covariance.
Candidates with high uncertainty are down-weighted to reduce false closures.
In parallel, semantic-language adapters project visual tokens into language-aware map attributes, enabling semantic mapping that can support text-conditioned retrieval and higher-level scene reasoning.

\subsection{Summary}

MACSLAM provides:
\begin{itemize}
    \item adaptive local visual-inertial odometry through cross-modal uncertainty correlation learning,
    \item improved trajectory consistency through unified multi-frame visual-inertial optimization,
    \item more reliable loop closure under domain shift via uncertainty-aware global matching,
    \item reduced system complexity by using one DINO-based front-end for geometry, retrieval, and semantics,
    \item semantic mapping with language-aligned features on top of the same SLAM backbone.
\end{itemize}
These capabilities support the thesis goal of a full spatial perception system that is accurate, robust, and deployable with minimal additional human engineering.